\documentclass[12pt]{article}
\usepackage{float}
\usepackage{graphicx}
\usepackage{fontspec}
\usepackage{url}
\usepackage{apacite}
\setmainfont{Times New Roman}
\usepackage[letterpaper,top=2cm,bottom=2cm,left=2cm,right=2cm,marginparwidth=1.5cm]{geometry}

\title{Weekly assignment 8: Ethic and Paradoxes}
\author{Truong Khang Tobias Nguyen}
\date{}

\begin{document}

\maketitle

The phrase "Quis custodiet ipsos custodes?" is Latin and translates to "Who guards the guards themselves?", raising questions about power, trust and accountability. Even though the phrase originally deals with ensuring marital fidelity, it is now more used to question the ability of people in positions of power to control their actions, and thus, not misuse their power.
\newline
In the context of cybersecurity, I believe that this phrase is a relevant question as cybersecurity professionals are entrusted with protecting sensitive information, assets and controlling access to said access. However, the same people we entrust with this task can abuse their power to compromise such information and assets themselves.
\newline
Therefore, challenges that arise is how to ensure that the people who protect us are held accountable themselves, how we can potentially regulate them or if we should regulate them in the first place.
\newline
As this dilemma (like most ethical questions) are not black-and-white problems and have more nuances, I argue that there is no solution to this challenge, but the most important thing is how we handle it.

\section{The Paradox}
I believe that the paradox is in the sentence itself. If an instance is responsible for overseeing others, who ensures that the overseer acts correctly. If another instance is appointed to oversee the previous instance, the same question arises again and subsequently, creating an infinitive loop of overseers, which is impossible to implement in practice.
\newline
Therefore, at the core of the paradox is a conflict between trust and control. Organizations need to trust their security professionals because they require access to sensitive systems and information to perform their work. At the same time, giving these professionals extensive power creates opportunities for abuse. If their own actions are not supervised, they can exploit their power while supervising them can undermine their efficiency and in a way, their professional autonomy. 

\section{Key Issues and Relevance for Cybersecurity}
The paradox is an ethical issue because the responsible use of power and the responsibilities that come with it are at the heart of the paradox itself. People entrusted with power have the duty to use it for the good of the people rather than for personal gain. In our context, this includes respecting confidentiality, protecting assets and information, and avoiding unauthorized access or surveillance.
\newline
For example, a system administrator may have access to employees' information and sensitive systems. Although this access is required for maintaining security, it also creates the responsiblity to avoid accessing information without reason. If the system administrator uses their power to surveil employees or use their personal information, they violate the trust and privacy of said employees.
\newline
This demonstrates a conflict between security and rights. Although monitoring is necessary to detect attacks, its excessive use can become an invasion of privacy - to paraphrase a former colleague from work, "Even though our information can't be accessed by everyone, there will always be one who can access it". The challenge here is that cybersecurity professionals must balance between security and rights rather than justifying every measure in the name of security.
\newline
Additionally, accountability is also an issue in this paradox. If cybersecurity professionals can investigate other while being beyond investigation themselves, they effectively possess unchecked power. That means that there is the need that their actions have to follow appropriate rules, that misuse has consequences, and that affected individuals can challenge abuse.
\newline
On that note, cybersecurity professionals usually operate with elevated privileges, allowing them to access systems and informations that ordinary users cannot. Even though this access is needed to protect organizations, it also leaves them as a source of risk.
\newline
For example, a malicious employee could steal employee or customer information, disable security controls, or manipulate logs to hide their activities. In that sense, the issue lies between giving cybersecurity professionals enough resources to operate while trusting them to not use said resources to harm the organization.

\section{Solving the Paradox?}
I do not believe that this paradox can be completly solved, but its risks can be reduced by distributing accountability. Rather than only having one single trusted instance, organizations should establish multiple sources of oversight.
\newline
First, cybersecurity professionals should receive only the resources needed to perform their work, similar to the least privilege principle. For particularly sensitive operations, it can require multiple people to approve an action which reduces the risk of an individual abusing their power.
\newline
Second, organizations should have logs that record privileged access. These logs should be audited by an independent instance rather than the ones whose actions are monitored. This ensures that organizations can detect whenever someone is using privileged access even if it is not required.
\newline
Last, organizations should establish clear standards of ethical conduct and consequences for misconduct. This creates a baseline on how cybersecurity professionals must act and possibly deter them from misuing their power.

\section{Pros and Cons}
The main advantage of such an approach that it reduces individual power. Multiple controls make it more difficult for individuals to misues their privileges, while logs ensure that misconduct can be discovered. This approach can also strenghten trust by showing that even cybersecurity professionals are subject to the principles that they enforce.
\newline
However, this approach also has disadvantages. Additional approval processes and audits increase costs, slow down decision-making, create administrative burden, and overall, become a bottleneck for organizations. Especially in emergencies, requiring multiple people to approve an action could delay a response to an attack.

\section{Conclusion}
The phrase introduces a fundamental problem on authority. Those tasked with protecting others can also use their power to harm them. In cybersecurity, this comes in the form of professionals with extensive accesses and privileges, introducing a technical side to this paradox as it combines security, privacy, trust, and responsibility.
\newline
My proposed solution is to distribute accountability through least privileges, logging privileged actions, and ethical standards. Although they don't solve the paradox, they can make abuse more difficult to perform and easier to detect. 
\newline
Ultimately, it is important that no instance should be considered beyond scrutiny and that responsible behaviour and accountability should be the norm for cybersecurity.

\end{document} 